% Options for packages loaded elsewhere
\PassOptionsToPackage{unicode}{hyperref}
\PassOptionsToPackage{hyphens}{url}
\PassOptionsToPackage{dvipsnames,svgnames,x11names}{xcolor}
\documentclass[
  10pt,
  fontset=fandol]{ctexart}
\usepackage{xcolor}
\usepackage[margin=16mm]{geometry}
\usepackage{amsmath,amssymb}
\setcounter{secnumdepth}{-\maxdimen} % remove section numbering
\usepackage{iftex}
\ifPDFTeX
  \usepackage[T1]{fontenc}
  \usepackage[utf8]{inputenc}
  \usepackage{textcomp} % provide euro and other symbols
\else % if luatex or xetex
  \usepackage{unicode-math} % this also loads fontspec
  \defaultfontfeatures{Scale=MatchLowercase}
  \defaultfontfeatures[\rmfamily]{Ligatures=TeX,Scale=1}
\fi
\usepackage{lmodern}
\ifPDFTeX\else
  % xetex/luatex font selection
\fi
% Use upquote if available, for straight quotes in verbatim environments
\IfFileExists{upquote.sty}{\usepackage{upquote}}{}
\IfFileExists{microtype.sty}{% use microtype if available
  \usepackage[]{microtype}
  \UseMicrotypeSet[protrusion]{basicmath} % disable protrusion for tt fonts
}{}
\makeatletter
\@ifundefined{KOMAClassName}{% if non-KOMA class
  \IfFileExists{parskip.sty}{%
    \usepackage{parskip}
  }{% else
    \setlength{\parindent}{0pt}
    \setlength{\parskip}{6pt plus 2pt minus 1pt}}
}{% if KOMA class
  \KOMAoptions{parskip=half}}
\makeatother
\usepackage{longtable,booktabs,array}
\newcounter{none} % for unnumbered tables
\usepackage{calc} % for calculating minipage widths
% Correct order of tables after \paragraph or \subparagraph
\usepackage{etoolbox}
\makeatletter
\patchcmd\longtable{\par}{\if@noskipsec\mbox{}\fi\par}{}{}
\makeatother
% Allow footnotes in longtable head/foot
\IfFileExists{footnotehyper.sty}{\usepackage{footnotehyper}}{\usepackage{footnote}}
\makesavenoteenv{longtable}
\setlength{\emergencystretch}{3em} % prevent overfull lines
\providecommand{\tightlist}{%
  \setlength{\itemsep}{0pt}\setlength{\parskip}{0pt}}
\usepackage{amsmath,amssymb,mathtools,needspace}
\xeCJKDeclareCharClass{CJK}{"2032,"0400->"04FF,"2160->"216B,"2460->"2473,"25A0,"25C7,"25CB,"25CF}
\IfFontExistsTF{Arial Unicode MS}{\xeCJKsetup{AutoFallBack=true}\setCJKfallbackfamilyfont{\CJKrmdefault}{Arial Unicode MS}}{}
\setcounter{secnumdepth}{0}
\setlength{\emergencystretch}{2em}
\allowdisplaybreaks
\usepackage{bookmark}
\IfFileExists{xurl.sty}{\usepackage{xurl}}{} % add URL line breaks if available
\urlstyle{same}
\hypersetup{
  pdftitle={差之毫厘{,}失之千里},
  colorlinks=true,
  linkcolor={Maroon},
  filecolor={Maroon},
  citecolor={Blue},
  urlcolor={Blue},
  pdfcreator={LaTeX via pandoc}}

\title{差之毫厘,失之千里}
\author{}
\date{}

\begin{document}
\maketitle

\subsubsection{12.3.4　差之毫厘，失之千里}\label{ux5deeux4e4bux6bebux5398ux5931ux4e4bux5343ux91cc}

加速因子究竟怎样选取才算合适呢？循着刘徽割圆术的思路，精确地计算直

\href{../../source-images/pdf-301.jpeg}{原页图像}

到正 3 072 边形的面积，直接选取加速因子 \(\omega=\dfrac13\)，再按式（12.3.3）计算，计算结果列于表 12.2 中，表中括弧 \(\langle\,\cdot\,\rangle\) 标明数据准确到小数点后第几位。\(\pi\) 的真值为 3.141 592 653 5\ldots。

\textbf{表 12.2}

{\def\LTcaptype{none} % do not increment counter
\begin{longtable}[]{@{}
  >{\raggedleft\arraybackslash}p{(\linewidth - 4\tabcolsep) * \real{0.3333}}
  >{\raggedleft\arraybackslash}p{(\linewidth - 4\tabcolsep) * \real{0.3333}}
  >{\raggedleft\arraybackslash}p{(\linewidth - 4\tabcolsep) * \real{0.3333}}@{}}
\toprule\noalign{}
\begin{minipage}[b]{\linewidth}\raggedleft
\(n\)
\end{minipage} & \begin{minipage}[b]{\linewidth}\raggedleft
\(S_n\)
\end{minipage} & \begin{minipage}[b]{\linewidth}\raggedleft
\(\widehat S_n\)
\end{minipage} \\
\midrule\noalign{}
\endhead
\bottomrule\noalign{}
\endlastfoot
12 & \(3.000\,000\,000\quad\langle0\rangle\) & \\
24 & \(3.105\,828\,541\quad\langle1\rangle\) & \(3.141\,104\,722\quad\langle3\rangle\) \\
48 & \(3.132\,628\,613\quad\langle1\rangle\) & \(3.141\,561\,971\quad\langle4\rangle\) \\
96 & \(3.139\,350\,203\quad\langle1\rangle\) & \(3.141\,590\,733\quad\langle5\rangle\) \\
192 & \(3.141\,031\,951\quad\langle3\rangle\) & \(3.141\,592\,534\quad\langle6\rangle\) \\
384 & \(3.141\,452\,472\quad\langle3\rangle\) & \(3.141\,592\,646\quad\langle7\rangle\) \\
768 & \(3.141\,557\,608\quad\langle4\rangle\) & \(3.141\,592\,653\quad\langle8\rangle\) \\
1 536 & \(3.141\,583\,892\quad\langle4\rangle\) & \(3.141\,592\,654\quad\langle8\rangle\) \\
3 072 & \(3.141\,590\,463\quad\langle5\rangle\) & \(3.141\,592\,654\quad\langle8\rangle\) \\
\end{longtable}
}

在割圆计算中，刘徽已获知直到正 3 072 边形的数据。以上数据表显示，利用这些数据按照刘徽加速技术进行精加工，即可获得祖冲之的``密率''3.141 592 65。

刘徽如果这样做，那么中华数学史乃至世界数学史上有关圆周率科学计算部分就要彻底改写了。

\textbf{差之毫厘，失之千里。刘徽的校正技术极为精彩，只是由于校正因子的处理稍欠精细，从而失之交臂，把一项千年称雄的数学成就留给了两百年后的祖冲之。}

\begin{center}\rule{0.5\linewidth}{0.5pt}\end{center}

\subsection{相关原页校注（agent 补充）}\label{ux76f8ux5173ux539fux9875ux6821ux6ce8agent-ux8865ux5145}

以下校注来自本节覆盖的原页，可能也涉及同页相邻小节；原文照录与校正说明分开保留。

\subsubsection{原书 PDF 页 301 的校注}\label{ux539fux4e66-pdf-ux9875-301-ux7684ux6821ux6ce8}

\href{../pages/pdf-301.md}{完整原页转录} · \href{../../source-images/pdf-301.jpeg}{原页图像}

\begin{quote}
校注（agent 补充）：表 12.2 的 \(S_n\) 数值从 3.000 000 000 开始，按原表保留；此前半径 \(r=10\) 寸的表 12.1 数值从 300 开始。本页未明说此处改用单位圆或作了归一化，转录未擅自改变尺度。
\end{quote}

\end{document}
